% concatenated from NN-eigen-rev.tex
\documentclass[reprint,superscriptaddress,aps,prl]{revtex4-2}
%\documentclass[preprint,superscriptaddress,aps,prl]{revtex4-2}


\usepackage[usenames,dvipsnames]{xcolor}
%\usepackage{xcolor}
\usepackage[utf8]{inputenc} 
\usepackage{bm}
%\usepackage{graphicx,psfrag}
\usepackage{graphicx}
\usepackage{booktabs}
\usepackage{amsmath,amsthm,amssymb,amsfonts}
\usepackage{mathtools}
\usepackage{physics}
%\usepackage{calc}
%\usepackage{makecell}
\usepackage{enumitem}
\usepackage{float}
%\usepackage{natbib}
\usepackage{tikz}
%\usepackage{soul,xcolor}
\usepackage[colorlinks=true,linkcolor=Blue,citecolor=Blue,urlcolor=Blue]{hyperref}

%\usepackage{xfp}

\usepackage[caption=false]{subfig}
\usepackage[capitalise]{cleveref}



\newcommand{\YZ}{\color{blue}}
\renewcommand{\P}{\mathbb{P}}
\newcommand{\R}{\mathbb{R}}
\newcommand{\N}{\mathbb{N}}

\newtheorem{teo}{Theorem}


\begin{document}

\title{A PINNs approach for the computation of eigenvalues in elliptic problems}


\author{Juli\'an Fern\'andez Bonder}
\email{jfbonder@dm.uba.ar}
\affiliation{Departamento de Matem\'atica, Facultad de Ciencias Exactas y Naturales, Universidad de Buenos Aires}
\affiliation{Instituto de C\'alculo UBA-CONICET, Buenos Aires, Argentina}

\author{Ariel Salort}
\email{ariel.salort@ceu.es}
\affiliation{Departamento de Matematicas y Ciencia de Datos, Universidad San Pablo-CEU, CEU Universities, Urbanizaci\'on Montepr\'incipe, 28660 Boadilla del Monte, Madrid, Spain}



\begin{abstract} 
In this paper, we propose a method for computing eigenvalues of elliptic problems using Deep Learning techniques. A key feature of our approach is that it is independent of the space dimension and can compute arbitrary eigenvalues without requiring the prior computation of lower ones. Moreover, the method can be easily adapted to handle nonlinear eigenvalue problems.
\end{abstract}

\maketitle

\section{Introduction}

The impact of artificial intelligence (AI) on everyday life has become undeniable in recent years. From smartphones to image recognition, from medical diagnostics to self-driving cars, machine learning—and deep learning in particular—has transformed the way we interact with technology.

This influence has also extended into mathematics and physics, where machine learning techniques are being applied to analyze complex phenomena that were traditionally approached using classical analytical methods. Among these, problems governed by partial differential equations (PDEs) are especially prominent, due to their central role in modeling physical systems across many disciplines.

Reflecting this trend, a growing number of software libraries and frameworks have been developed as educational and research tools for solving PDE-based problems in computational science and engineering; see, for instance, \cite{Lu-Meng-Mao, Rackauckas, Rackauckas-Innes-Ma-Bettencourt-etal} and references therein.

In parallel, neural network-based approaches for approximating solutions to PDEs have gained increasing attention. These include successful applications to problems such as the Burgers, Eikonal, and heat equations \cite{Blechschmidt-Ernst}, as well as linear diffusion equations in complex two-dimensional geometries \cite{Berg-Nystrom}.

A natural extension of this idea is the use of neural networks (NNs) as general-purpose function approximators for solving differential equations. A prototypical example arises in quantum mechanics, where one seeks to solve the eigenvalue problem
\begin{equation}\label{eigen}
-\Delta u + V(x) u = E u,\quad \text{in } \R^n,
\end{equation}
with a given potential function \( V \colon \R^n \to \R \) and eigenvalue parameter \( E \in \R \). The idea of employing NNs to approximate solutions of such equations dates back to the 1990s \cite{Lagaris-Likas-Fotiadis}, but the field gained significant momentum with the introduction of {\em Physics-Informed Neural Networks} (PINNs) in \cite{Raissi-Perdiakaris-Karniadakis}, which triggered a wave of research and applications (see also \cite{Raissi-Perdiakaris-Karniadakis2, Karniadakis-etal}).

Most early PINN-based work focused on source problems, where the term \( E u \) in \eqref{eigen} is replaced by a known forcing function \( f \). The eigenvalue problem itself has been more recently addressed: first in one dimension \cite{Jin-Mattheakis-Protopapas1, Jin-Mattheakis-Protopapas2, Sakar}, and later in two dimensions \cite{Holliday-Lindner-Ditto}. Related approaches have also been proposed in \cite{Rowan-Evans-Maute-Dootsan}, which include extensions to nonlinear eigenvalue problems.

 
Several methodologies have been proposed to address eigenvalue problems using neural networks.
The approach in \cite{Sakar} constructs a trial wavefunction that automatically satisfies the boundary conditions, represents the unknown component with a neural network, and trains the network by minimizing an appropriately defined loss function.
%
In \cite{Jin-Mattheakis-Protopapas1, Jin-Mattheakis-Protopapas2}, the authors introduce an unsupervised neural network architecture that simultaneously learns both the eigenvalues and the corresponding eigenfunctions through a scanning mechanism.
%
Alternatively, \cite{Rowan-Evans-Maute-Dootsan} employs a neural network to discretize the eigenfunction and minimize the associated Rayleigh quotient.

\normalcolor



\begin{figure}[tb]
\centering
\includegraphics[width=.9\columnwidth]{NN-Tikz.png}
\caption{A fully connected feedforward NN with an \(n\)-dimensional input and a 1-dimensional output.}
\label{NN.fig}
\end{figure}

In these works, the solution \( u \) is approximated using a fully connected feedforward neural network (see Fig.~\ref{NN.fig}), with variations in the number of hidden layers and activation functions across implementations. To approximate the eigenvalue \( E \), the authors minimize the {\em Rayleigh quotient}:
\[
L(u) = \frac{\int |\nabla u|^2 + V(x) u^2\, dx}{\int u^2\, dx}.
\]
Minimizing \( L(u) \) yields an approximation of the principal eigenvalue \( E_1 \) and its associated eigenfunction \( u_1 \). Higher-order eigenpairs can be computed by iteratively minimizing the same functional while constraining the search to subspaces orthogonal to previously computed eigenfunctions.

This approach, while effective, presents two main limitations:
\begin{enumerate}
\item It is inherently sequential: to compute eigenvalues within a given range, one must first compute all lower ones, which can be inefficient.
\item Its extension to nonlinear problems (such as those involving the \(p\)-Laplace operator) is limited, since orthogonality conditions between eigenfunctions are no longer available.
\end{enumerate}


%\section{Introduction}
%
%The impact that AI methods have had on our everyday lives in recent years is now undeniable. From smartphones to image recognition, from medical diagnostics to self-driving cars, machine learning (of which deep learning is a subset) has fundamentally changed the way we live today.
%
%In particular, in mathematics and physics, many phenomena that were traditionally analyzed using classical analytical tools are now being studied with machine learning techniques, leading to new and often more effective insights into longstanding problems.
%
%In recent years, a growing number of software libraries and frameworks have been developed as educational and research tools for addressing problems in computational science and engineering—particularly those involving partial differential equations (PDEs) and their diverse applications. Notable examples include \cite{Lu-Meng-Mao, Rackauckas, Rackauckas-Innes-Ma-Bettencourt-etal} and references therein.
%
%In parallel, alternative approaches based on neural networks have emerged for approximating solutions to PDEs such as the Burgers, Eikonal, and heat equations; see \cite{Blechschmidt-Ernst} for a recent overview. In particular, \cite{Berg-Nystrom} introduces an effective neural network method for linear diffusion problems posed in complex two-dimensional domains.
%
%As the previous discussion exemplifies, the use of Neural Networks (NNs) as function approximators has become a standard technique. In physics, for instance, many problems can be formulated as the task of solving a differential equation (DE). One of the most well-known examples comes from quantum mechanics and involves solving the so-called {\em eigenvalue problem}
%\begin{equation}\label{eigen}
%-\Delta u + V(x) u = E u,\quad \text{in } \R^n
%\end{equation}
%where $V\colon \R^n\to\R$ is the potential function and $E\in\R$ is the eigenvalue parameter. The use of NNs to solve DEs such as this one dates back to \cite{Lagaris-Likas-Fotiadis} in the 1990s, but it was not until the late 2010s that the term {\em Physics-Informed Neural Networks} (PINNs) was introduced in \cite{Raissi-Perdiakaris-Karniadakis}, sparking an explosion of research in this area (see also \cite{Raissi-Perdiakaris-Karniadakis2, Karniadakis-etal} and further references therein).
%
%
%
%Recall that the aforementioned papers primarily addressed the source problem, i.e. where the term $E u$ in \eqref{eigen} is replaced by a known forcing term $f$. The eigenvalue problem \eqref{eigen} itself was recently studied in the one-dimensional case in \cite{Jin-Mattheakis-Protopapas1, Jin-Mattheakis-Protopapas2} and later extended to the two-dimensional case in \cite{Holliday-Lindner-Ditto}.  See also \cite{Rowan-Evans-Maute-Dootsan} where the same problem is tackled using similar methods and an extension to the nonlinear eigenvalue problem is also discussed.
%
%\begin{figure}[tb]
%\centering
%\includegraphics[width=.9\columnwidth]{NN-Tikz.png}
%\caption{A fully connected feedforward NN with and $n-$dimensional input and a $1-$dimensional output.}
%\label{NN.fig}
%\end{figure}
%
%In these works, the authors used a fully connected feedforward neural network (NN) to approximate the function $u$ that solves \eqref{eigen}. The activation function, as well as the number of hidden layers, vary from one article to the other (see Fig.~\ref{NN.fig}). To compute the eigenvalues using this NN the authors proposed minimizing the {\em Rayleigh quotient}:
%$$
%L(u) = \frac{\int |\nabla u|^2 + V(x) u^2\, dx}{\int u^2\, dx}
%$$
%The network is trained by minimizing $L(u)$ yielding an approximation of $E_1$, the lowest (or {\em principal}) eigenvalue of \eqref{eigen}, along with the corresponding eigenfunction $u_1$. To compute higher order eigenvalues and eigenfunctions, the same loss functional $L(u)$ is minimized again, but this time restricted to the subspace orthogonal to $u_1$. Repeating this procedure allows one to compute successive eigenpairs.
%
%This method presents two main drawbacks: 
%\begin{enumerate}
%\item The approach is inherently sequential. If one wishes to compute eigenvalues within a given interval, say $[E_*, E^*]$, it is necessary to compute all preceding eigenvalues starting from the first, which can be computationally inefficient.
%
%\item Extending this method to nonlinear problems (such as those involving the $p-$Laplace operator) is limited to the computation of the first eigenvalue, since in such cases there is no natural orthogonality condition between eigenfunctions.
%\end{enumerate}

\section{Our approach}

To overcome the drawbacks mentioned above, we propose a different approach to problem \eqref{eigen}. We define the following loss function:
\begin{equation}\label{Lambda}
\Lambda(u, E) = \frac{\int (\Delta u - V(x) u + Eu)^2\, dx}{\int u^2\, dx}.
\end{equation}
Given a fixed $E\in \R$, minimizing with respect to $u$ yields a {\em loss curve}:
\begin{equation}\label{loss-curve}
\Lambda(E) = \inf_u \Lambda(u, E).
\end{equation}
A proof of the existence of a minimizer $u_E$ for \eqref{loss-curve} is provided in Theorem \ref{teo.min} in the appendix.

It is relatively straightforward to derive a theoretical upper bound for the loss curve $\Lambda(E)$. Indeed, let $u_k$ denote the $k-$th eigenfunction to \eqref{eigen}, normalized such that $\int u_k^2\ dx=1$. Then we have:
$$
\Lambda(E)\le \int (\Delta u_k - V(x) u_k + Eu_k)^2\, dx = (E_k-E)^2,
$$
and consequently,
$$
\Lambda(E) \le \min_{k} (E_k-E)^2.
$$
See Fig.~\ref{Lambda.fig}, where this upper bound is plotted in the two-dimensional case with a confining potential given by $V(x)=0$ for $|x|\le 1$ and $V(x)=\infty$ for $|x|>1$. In this case the eigenvalues are known explicitly. The first four are:
$$
 E_1 = 5.7832,\ E_2 = 14.6819,\ E_3 =  26.3743,\ E_4 =  30.4713.
 $$
\begin{figure}[tb]
\centering
\includegraphics[width=.9\columnwidth]{Loss.png}
\caption{The loss curve $\Lambda(E)$ in $2-$dimensional space with a confined potential in the unit disk.}
\label{Lambda.fig}
\end{figure}
Our approach proceeds as follows:

We first fix the interval $[E_*, E^*]$  where the eigenvalues are expected to lie. Then, we take a uniform partition of this interval:
$$
E_*=E^1<E^2<\cdots<E^j=E^*,
$$ 
and for each $E^i$ we train our neural network to minimize $\Lambda(u, E^i)$, obtaining a corresponding minimizer $u^i$.

Finally, we identify the minima of the loss curve $\Lambda(E^i)$. The values $E^i$ at which these minima occur approximate the eigenvalues of \eqref{eigen}, and the associated minimizers $u^i$ approximate the corresponding eigenfunctions.

 
The aim of this article is to give an approach to compute eigenvalues  based on neural networks whose implementation is simple and no need of previous eigenvalues is required.
Our goal is not to obtain the best estimates for lineal problems (for that aim there is efficient and well-stuided methods based on finite differences/elements).  However, our method is highly flexible and has several adventages over other method of the literature as explained as follows.


The aim of this article is to introduce a neural network–based approach for computing eigenvalues, which  is straightforward to implement and does not require any prior knowledge of the eigenvalues. We emphasize that our objective is not to achieve the most accurate estimates for linear problems, as efficient and well-established methods based on finite differences or finite elements already exist for that purpose. However, the proposed method offers significant flexibility and several advantages over existing approaches in the literature, as detailed below.


Advantages of our method:
\begin{enumerate}
\item  Unlike the aforementioned methods, our approach allows the computation of higher eigenvalues without requiring the prior computation of lower ones.


\item While other methods often face challenges when computing eigenvalues in higher dimensions, our approach has been shown to be effective regardless of the dimensionality.

\item The computation of eigenvalues is robust and works effectively independent of the structure of the potential function $V$ or the geometry of the domain under consideration.

\item The neural network architecture employed in our implementation is simple and extends naturally to nonlinear cases. This enables the computation of higher eigenvalues, producing novel results in several contexts shown ir our examples. While we do not claim more accuracy compared to existing methods, our approach provides the computation of higher eigenvalues even in higher-dimensional settings.

 
\item Finally, our approach can be readily adapted to more general operators beyond the $p$-Laplacian.
\end{enumerate}


\normalcolor

\section{Implementation details}

We consider a fully connected feedforward neural network (as in Fig.~\ref{NN.fig}) with two hidden layers and the hyperbolic tangent activation function, $\tanh$. 

 


We consider a  fully connected neural network,  implemented as a feed-forward multilayer perceptron with three hidden layers. Each hidden layer contains 128 neurons with $\tanh$ activation functions. The network maps a  spatial coordinate $x\in \R^d$ to a scalar output and is defined as
$$
x \mapsto \text{NN}(x) \in \mathbb{R}
$$
All linear layers are initialized using Xavier uniform initialization with the appropriate gain for the tanh activation, and all biases are set to zero.

To enforce the homogeneous Dirichlet boundary condition on the domain $\Omega$, the raw network output $\tilde u(x)$ is multiplied by a boundary distance function $B(x)$, chosen so that
$$
B(x)=0 \text{ for all } x\in\partial\Omega, \qquad B(x)>0 \text{ for all } x\in\Omega.
$$
This construction ensures that the final approximation automatically vanishes on the boundary. Thus, the network returns
$$
u(x)=B(x)\tilde  u(x),
$$
where $\tilde  u(x)$ is the output of the multilayer perceptron.
\normalcolor



 
We begin by studying the simpler case, where confining potentials, as described in the previous section, are considered. So, we let
\normalcolor
\begin{equation}\label{conf.V}
V(x) = \begin{cases}
0 & \text{if } x\in \Omega\\
\infty & \text{if } x\not\in \Omega,
\end{cases}
\end{equation}
which is equivalent to solving the Helmholtz problem with homogeneous Dirichlet boundary condition:
\begin{equation}\label{eigen2}
\begin{cases} 
-\Delta u = E u & \text{in }\Omega\\
u=0 & \text{on }\partial\Omega.
\end{cases}
\end{equation}

In this case, to  impose the boundary condition, we multiply the output of the NN by a boundary factor $B(x)=1-x^2$ that is a smooth approximation to the distance to the boundary. To ensure normalization, we include a penalization term in the loss function to enforce $\|u\|_2=1$. The resulting loss is:
\begin{equation}\label{lossmu}
\bar \Lambda(u, E) = \int_\Omega (\Delta u + Eu)^2\, dx + \mu\left(\int_\Omega u^2\, dx - 1\right)^2.
\end{equation}
All integrals are approximated using a Monte Carlo sampling scheme. 

As described in the previous section, we train the NN to minimize $\bar\Lambda(u, E)$ for each value in a discrete partition $E_* = E^1<\cdots<E^j=E^*$. To improve efficiency, for each $E^i$, the network is initialized with the weights and biases obtained from training at $E^{i-1}$.

After training is complete, we obtain the discrete loss curve
$$
\bar \Lambda(E^i)=\bar\Lambda(u^i, E^i)
$$
and each local minimum below a given threshold $\epsilon>0$ is taken as an approximate eigenvalue for \eqref{eigen2}. 
\begin{figure}[tb]
\centering
\includegraphics[width=.9\columnwidth]{lossNN-lento.png}
\caption{The loss curve $\bar\Lambda(E)$ in two-dimensional space with a confined potential in the unit disk compared with the theoretical upper bound.}
\label{lossNN.fig}
\end{figure}

Fig.~\ref{lossNN.fig} shows the loss curve $\bar\Lambda(E^i)$ in the case where $\Omega$ is the unit disk in $\R^2$. The computed eigenvalues are accurate up to the resolution of the discretized $E^i-$grid. Naturally, once an approximate eigenvalue $E^{i_0}$ is detected, the method can be refined locally by using a finer grid in the interval $[E^{i_0-1}, E^{i_0+1}]$ to improve precision.

The associated computed eigenfunctions are displayed in Fig.~\ref{func.fig}
\begin{figure}[tb]
\centering
\begin{tabular}{cc}
\includegraphics[width=.4\columnwidth]{eigen1-1.png} & \includegraphics[width=.4\columnwidth]{eigen1-2.png} \\
\includegraphics[width=.4\columnwidth]{eigen2-1.png} & \includegraphics[width=.4\columnwidth]{eigen2-2.png} \\
\includegraphics[width=.4\columnwidth]{eigen3-1.png} & \includegraphics[width=.4\columnwidth]{eigen3-2.png} \\
\includegraphics[width=.4\columnwidth]{eigen4-1.png} & \includegraphics[width=.4\columnwidth]{eigen4-2.png} \\
\end{tabular}
\caption{The first four computed eigenfunctions of \eqref{eigen2} in the $2-$dimensional unit disk.}
\label{func.fig}
\end{figure}

\section{Further examples}

In this section, we illustrate the flexibility and robustness of our method by applying it to a variety of settings beyond the two-dimensional unit disk with a confining potential. These examples include the computation of higher-order eigenvalues, problems in higher spatial dimensions, domains with different geometries, and more general potential functions.


All computations and plots presented in this article were generated using the code provided in the ‘Code Availability’ section.
\normalcolor

\subsection{Higher eigenvalues}
Our method enables the computation of higher-order eigenvalues without relying on sequential orthogonalization procedures. By choosing an appropriate interval $[E_*, E^*]$ and refining the discretization grid, we can directly recover eigenvalues that are far from the principal one.

One technical point that must be considered when computing higher eigenvalues is the appropriate choice of the hyperparameter $\mu$ in the loss function \eqref{lossmu}. For a randomly initialized NN output $u$, the first term in the loss is typically of the order
$$
\int_\Omega (\Delta u + E u)^2\, dx \sim c_1 E^2 + c_2,
$$
while the normalization penalty term is of the order
$$
\mu \left( \int_\Omega u^2\, dx - 1 \right)^2 \sim c_3 \mu.
$$
Hence, if $E^2 \gg \mu$, the loss is dominated by the first term, and the network tends to minimize it by driving $u \approx 0$, which is undesirable. To maintain a proper balance between the two terms and ensure meaningful training, it is necessary to scale $\mu \propto E^2 $.

To assess the accuracy of our method, we computed the eigenvalues of problem \eqref{eigen2} in the two-dimensional unit square $[0,1]\times[0,1]$ that lie within the interval $[44, 55]$. In this range, problem \eqref{eigen2} has one eigenvalue, denoted $E_{(1)}$, with value $E_{(1)} = 49.348$. See Fig.~\ref{eigen55.fig} for the computed loss curve in this interval, together with the theoretical upper bound.



%\medskip
%\underline{Parameters}: \\
%$\Omega=[0,1]\times[0,1]$\\
%V=0\\
%Sample points: 5000\\
%mu=500\\
%n\_epochs=5000\\
%early\_stopping\_patience=300\\
%E\_min=0\\
%E\_max=105\\
%E\_step=80

 
For illustrative purposes, we compute the eigenvalues of \eqref{eigen2} in the two-dimensional unit square.   In Table \ref{table.1} we compare the exact eigenvalues with those obtained by our method, together with the relative error   and   residuals. In Figure \ref{loss.lap.square.fig} we display the loss curve $\tilde \Lambda (E)$ alongside the theoretical upper bound, as well as the first two eigenfunctions.

These computations can be readily carried out for different potentials; however, in such cases, exact eigenvalues are not available for comparison.


\begin{table*}[h!] 
\centering 

\begin{tabular}{c c|c|c|c|c}
$m$ & $n$ & $E_{(m,n)} = \pi^2 (m^2 + n^2)$ & Our method & Relative error (\%)&  Residual \\ 
\hline
1 & 1 & $2\pi^2$   $\approx 19.739$ & 19.739033 & $8.9\times 10^{-6}$ & $6.72\times 10^{-3}$ \\
2 & 1 & $5\pi^2$   $\approx 49.348$ & 49.385078 & $7.5\times 10^{-4}$ & $1.16\times 10^{-1}$\\
1 & 2 & $5\pi^2$   $\approx 49.348$ & 49.385078 & $7.5\times 10^{-4}$ & $1.16\times 10^{-1}$\\
2 & 2 & $8\pi^2$   $\approx 78.956$  & 78.962188 & $1.05\times 10^{-5}$& $1.26\times 10^{-1}$  \\
3 & 1 & $10\pi^2$  $\approx 98.696$ & 98.594123 & $1.1\times 10^{-3}$ & $1.40\times 10^{-4}$\\
1 & 3 & $10\pi^2$  $\approx 98.696$ & 98.594123 & $1.1\times 10^{-3}$ & $1.40\times 10^{-4}$
\end{tabular}
\caption{Eigenvalues of problem \eqref{eigen2} in the two-dimensional unit square $[0,1]\times[0,1]$ compared with the exact values.}
\label{table.1}

\end{table*}

\normalcolor

\begin{figure}[tb]
\centering
\includegraphics[width=.9\columnwidth]{loss_2d_cuadrado_entre_45_y_55_1er_autovalor.png}
\caption{The loss curve $\bar\Lambda(E)$ in two-dimensional space  in the unit square computed in the interval [44,55].}
\label{eigen55.fig}
\end{figure}


\subsection{Higher dimensions}

The proposed framework extends naturally to higher spatial dimensions, with minimal changes required in the implementation. As an example, we consider the unit ball in $\R^d$, with $d=3,4$ where the eigenvalue problem becomes:
$$
\begin{cases}
-\Delta u = Eu & \text{in } B^d\\
u=0 & \text{on }\partial B^d.
\end{cases}
$$
Despite the increased computational cost, the method remains effective, as shown in Fig.~\ref{loss3D.fig}.




%\textbf{Example 7.} Laplacian with a potential in the unit ball in  $\mathbb{R}^3$
%
%\medskip
%\underline{Parameters}: \\
%$\Omega=B(0,1)$ in $\mathbb{R}^4$\\
%V=0\\
%Sample points: 5000\\
%mu=200\\
%n\_epochs=2000\\
%early\_stopping\_patience=300\\
%E\_min=0\\
%E\_max=55\\
%E\_step=1000
%

 
In dimension $d=3$, exact formulas for the eigenvalues are known, for $m,n\in \N$, the eigenvalues are computed as $E_{(m,n)}=j_{m,n}^2$, where $j_{m,n}$ denotes the $n-$th positive zero of the Bessel function $J_{m+\frac12}$. Similarly, when $d=4$, the eigenvalues are given by $E_{(m,n)}=y_{m,n}^2$, where $y_{m,n}$ denotes the $n-$th positive zero of the Bessel function $J_{m+1}$.

In Tables \ref{table.2} and \ref{table.3}, we compare the exact eigenvalues with those obtained by our method.

\normalcolor
\begin{table*}[h!]
\centering

\begin{tabular}{cc|c|c|c|c}
$m$ & $n$ & $E_{(m,n)} = j_{m,n}^2$ & Our method & Relative error (\%) & Residual \\  
\hline
0 & 1 & $j_{0,1}^2$   $\approx 9.8686$ & 9.885022 & $1.66\times 10^{-3}$& 8.61$\times 10^{-2}$ \\
1 & 1 & $j_{1,1}^2$   $\approx 20.1907$ & 20.182569 & $4.02\times 10^{-4}$ & 4.59$\times 10^{-2}$\\
2 & 1 & $j_{2,1}^2$   $\approx 33.2251$ & 33.243410 & $5.51\times 10^{-4}$& 1.00$\times 10^{-1}$\\
0 & 2 & $j_{0,2}^2$   $\approx 39.4784$  & 39.494155  & $3.99\times10^{-4} $& 1.81$\times 10^{-1}$ \\
3 & 1 & $j_{3,1}^2$   $\approx 48.9737$  & 48.826976  & $2.99\times 10^{-3}$& 2.24$\times 10^{-1}$ \\
1 & 2 & $j_{1,2}^2$   $\approx 59.6946$  & 59.734069   & $6.61\times 10^{-4}$ & 2.04$\times 10^{-1}$ \\
4 & 1 & $j_{4,1}^2$   $\approx 66.9489$  & 66.991853  & $6.41\times 10^{-4}$ & 4.82$\times 10^{-1}$ 
\end{tabular}
\caption{Eigenvalues of problem \eqref{eigen2} in the unit ball in $\R^3$ compared with the exact values. }
\label{table.2}
 
\end{table*}
 


%
%\textbf{Example 6.} Laplacian with a potential in the unit ball in  $\mathbb{R}^4$
%
%\medskip
%\underline{Parameters}: \\
%$\Omega=B(0,1)$ in $\mathbb{R}^4$\\
%V=0\\
%Sample points: 5000\\
%mu=200\\
%n\_epochs=2000\\
%early\_stopping\_patience=300\\
%E\_min=0\\
%E\_max=55\\
%E\_step=1000
%


\begin{table*}[h!]
\centering
 
\begin{tabular}{cc|c|c|c|c}
$m$ & $n$ & $E_{(m,n)} = y_{m,n}^2$ & Our method & Relative error (\%) & Residual \\ 
\hline
0 & 1 & $y_{0,1}^2$   $\approx 14.6819$ & 14.641796 & $2.73\times 10^{-3}$ & 3.84$\times 10^{-2}$ \\
1 & 1 & $y_{1,1}^2$   $\approx 26.3746$ & 26.257056 & $4.47\times 10^{-3}$& 2.13$\times 10^{-1}$\\
2 & 1 & $y_{2,1}^2$   $\approx 40.7064$ & 40.794700 & $2.16\times 10^{-3}$& 1.98$\times 10^{-1}$\\
0 & 2 & $y_{0,2}^2$   $\approx 49.2184$  & 49.175395  & $8.74\times 10^{-4}$& 6.04$\times 10^{-1}$ \\
\end{tabular}
\caption{Eigenvalues of problem \eqref{eigen2} in the unit ball in $\R^4$ compared wih the exact values. }
\label{table.3}

\end{table*}
  



%\begin{figure}[tb]
%\centering
%\includegraphics[width=.9\columnwidth]{Loss.png}
%\caption{The loss curve $\Lambda(E)$ in $2-$dimensional space with a confined potential in the unit disk.}
%\label{Lambda.fig}
%\end{figure}


\begin{figure}[tb]
\centering
\includegraphics[width=.9\columnwidth]{loss_n=3_V=0_p=2_mu=100_reseteo_cada_15_iteraciones.png}
\includegraphics[width=.9\columnwidth]{loss_n=4_p=2_V=0_mu=200.png}
\caption{The loss curve $\bar\Lambda(E)$ in three-dimensional space (above) and in the four-dimensional space (below) in the unit ball compared with the theoretical upper bound.}
\label{loss3D.fig}
\end{figure}




%\begin{figure}[tb]
%\centering
%\begin{tabular}{cc}
%\includegraphics[width=.4\columnwidth]{n=2 p=2 v=0 mu=150 cuadrado autof1.png} & \includegraphics[width=.4\columnwidth]{n=2 p=2 v=0 mu=150 cuadrado autof13d.png} \\
%\includegraphics[width=.4\columnwidth]{n=2 p=2 v=0 mu=150 cuadrado autof2.png} & \includegraphics[width=.4\columnwidth]{n=2 p=2 v=0 mu=150 cuadrado autof23d.png} \\
%\includegraphics[width=.4\columnwidth]{n=2 p=2 v=0 mu=150 cuadrado autof3.png} & \includegraphics[width=.4\columnwidth]{n=2 p=2 v=0 mu=150 cuadrado autof33d.png} \\
%\includegraphics[width=.4\columnwidth]{n=2 p=2 v=0 mu=150 cuadrado autof4.png} & \includegraphics[width=.4\columnwidth]{n=2 p=2 v=0 mu=150 cuadrado autof43d.png} \\
%\end{tabular}
%\caption{The first four eigenfunctions of \eqref{eigen2} in the $2-$dimensional unit square.}
%\label{func.square.fig}
%\end{figure}












\subsection{Different geometries}
We also explore domains with more complex geometries, such as squares, annuli, or triangles. The boundary factor $B(x)$ can be adapted to these shapes, ensuring that the Dirichlet condition is enforced. In Fig.~\ref{loss.lap.square.fig}, we show eigenfunctions obtained in a square, in Fig.~\ref{loss.lap.annulus.fig}  the annular region is considered and in Fig.~\ref{loss.lap.triang.fig} the results in a triangle are shown.  
In Table \ref{table.4},  the exact eigenvalues of \eqref{eigen2} in a triangle are compared with those obtained by our method. 

%\underline{Parameters}: \\
%$\Omega$=triangle with vertices $(0,0)$, $(0,1)$, $(1,0)$\\
%V=0\\
%Sample points: 5000\\
%mu=1500\\
%n\_epochs=2000\\
%early\_stopping\_patience=300\\
%E\_min=0\\
%E\_max=180\\
%E\_step=120

\begin{table*}[h!]
\centering
 
\begin{tabular}{c c|c|c|c|c}
$m$ & $n$ & $E_{m,n} = \pi^2 (m^2 + n^2)$ & Our method & Relative error (\%) & residual \\ 
\hline
1 & 2 & $5\pi^2$   $\approx 49.348$ & 49.369942 & $4.43\times 10^{-4}$ & 2.42$\times 10^{-1}$ \\
2 & 1 & $5\pi^2$   $\approx 49.348$ & 49.369942 & $4.43\times 10^{-4}$ & 2.42e$\times 10^{-1}$\\
1 & 3 & $10\pi^2$   $\approx 98.696$ & 98.892161 & $4.45\times 10^{-7}$ & 1.26e$\times 10^{0}$\\
3 & 1 & $10\pi^2$   $\approx 98.696$  & 98.892161 & $4.45\times 10^{-7}$ & 1.26$\times 10^{0}$  \\
2 & 3 & $13\pi^2$  $\approx 128.996$ & 128.048509 & $2.00\times 10^{-3}$ & 1.98$\times 10^{0}$\\
3 & 2 & $13\pi^2$  $\approx 128.996$ & 128.048509 & $2.00\times 10^{-3}$ & 1.98$\times 10^{0}$ \\
1 & 4 & $17\pi^2$  $\approx 167.551$ & 167.687185 & $5.73\times 10^{-4}$ & 2.16$\times 10^{0}$\\
4 & 1 & $17\pi^2$  $\approx 167.551$ & 167.687185 & $5.73\times 10^{-4}$ &  2.16$\times 10^{0}$ 
\end{tabular}
\caption{Eigenvalues of problem \eqref{eigen2} in the two-dimensional triangle with   vertices $(0,0)$, $(0,1)$, $(1,0)$ compared with the exact values.}
 
\label{table.4}
\end{table*}

\begin{figure}[tb]
\centering
\includegraphics[width=.9\columnwidth]{loss_n=2_p=2_V=0_mu=450_cuadrado.png}
\includegraphics[width=.45\columnwidth]{n=2_p=2_v=0_mu=150_cuadrado_autof1.png}
\includegraphics[width=.45\columnwidth]{n=2_p=2_v=0_mu=150_cuadrado_autof13d.png}
\includegraphics[width=.45\columnwidth]{n=2_p=2_v=0_mu=150_cuadrado_autof2.png}
\includegraphics[width=.45\columnwidth]{n=2_p=2_v=0_mu=150_cuadrado_autof23d.png}
\caption{The loss curve $\bar\Lambda(E)$ in two-dimensional space in the unit square compared with the theoretical upper bound. Below it is displayed the first two eigenfunction.}
\label{loss.lap.square.fig}
\end{figure}

\begin{figure}[tb]
\centering
\includegraphics[width=.9\columnwidth]{loss_n=2_p=2_V=0_annulus.png}
\includegraphics[width=.45\columnwidth]{eigen1-annulus.png}
\includegraphics[width=.45\columnwidth]{eigen1-3d-annulus.png}
\caption{The loss curve $\bar\Lambda(E)$ in two-dimensional space in the unit annulus of radii 0.5 and 1 compared with the theoretical upper bound. Below it is displayed the first eigenfunction.}
\label{loss.lap.annulus.fig}
\end{figure}

\begin{figure}[tb]
\centering
\includegraphics[width=.9\columnwidth]{loss_n=2_p=2_V=0_mu=15000_triangulo.png}
%\includegraphics[width=.45\columnwidth]{eigen1-triangulo.png}
%\includegraphics[width=.45\columnwidth]{eigen1-3d-triangulo.png}
%\includegraphics[width=.45\columnwidth]{eigen_p=2_n=2_V=0_triangulo_3.png}
%\includegraphics[width=.45\columnwidth]{eigen_p=2_n=2_V=0_triangulo_3_3d.png}
\includegraphics[width=.45\columnwidth]{eigen_p=2_n=2_V=0_triangulo_4.png}
\includegraphics[width=.45\columnwidth]{eigen_p=2_n=2_V=0_triangulo_4_3d.png}
\includegraphics[width=.45\columnwidth]{eigen_p=2_n=2_V=0_triangulo_5.png}
\includegraphics[width=.45\columnwidth]{eigen_p=2_n=2_V=0_triangulo_5_3d.png}
\caption{The loss curve $\bar\Lambda(E)$ in two-dimensional space  in the triangle of vertices $(0,0)$, $(0,1)$, $(1,0)$  compared with the theoretical upper bound and the third and  fourth computed eigenfunctions.}
\label{loss.lap.triang.fig}
\end{figure}



\subsection{Different potentials}
 
Beyond hard-wall confinement, our approach accommodates general potential functions $V(x)$, including smooth wells, multi-well potentials, step-like profiles, compactly supported wells, among others. 

As an example, we consider a harmonic potential $V(x)=\frac{\omega^2}{2}|x|^2$ for different values of $\omega$, illustrating the method's applicability to physically relevant models. Fig.~\ref{pot1.fig} displays loss curve and the first eigenfunction for $\omega=1$ and Fig.~\ref{pot2.fig} for $\omega=10$.

 
In Fig.~\ref{fig.well-two}, we display the three-dimensional double-well Gaussian potential $V(x)=-e^{-2|x-x_0|^2} - e^{-2|x+x_0|^2}$, $x_0=(1,0,0)$, together with the loss curve. 

Finally, in Fig.~\ref{fig.well-four} we consider a  three-dimensional four-well compactly supported  potential and we plot the corresponding loss curve. More precisely, given $w = (x,y,z) \in \mathbb{R}^3$, and define the centers of the four wells as $c_1 = (a,a,0)$, $c_2 = (-a,a,0)$, $c_3 = (a,-a,0)$, $c_4 = (-a,-a,0)$, where $a=0.22$. Also, let $r_i = \| w - c_i \|$, $\tilde r_i = \frac{r_i}{\sigma}$ with $\sigma = 0.2$. We define the compactly supported bump function
$$
\phi(\tilde r) =
\begin{cases}
\exp(-(1-\tilde r^2)^{-1}), & \tilde r < 1,\\
0, & \tilde r \ge 1.
\end{cases}
$$
Then, the potential $V(w)$ is defined as
$
V(x) = - \sum_{i=1}^{4} \phi\!\left( \frac{\| x - c_i \|}{\sigma} \right), \quad V_0 = 1.
$
%Explicitly, using the centers \((\pm a, \pm a, 0)\), the potential can be written as
%$$
%V(x,y,z) = - \sum_{(\epsilon_1, \epsilon_2) \in \{\pm 1\}^2} 
%\exp \Bigg( - \frac{1}{
%1 - \frac{(x - \epsilon_1 a)^2 + (y - \epsilon_2 a)^2 + z^2}{\sigma^2} }
%\Bigg) 
%\mathbf{1}_{\{ (x - \epsilon_1 a)^2 + (y - \epsilon_2 a)^2 + z^2 < \sigma^2 \} } .
%\]


 



\begin{figure}[tb]
\centering
\includegraphics[width=.8\columnwidth]{loss_n=2_V=harmonic_p=2.png}
\includegraphics[width=.45\columnwidth]{eigen_n=2_V=harmonic_p=2_1.png}
\includegraphics[width=.45\columnwidth]{eigen_n=2_V=harmonic_p=2_1_3d.png}
\caption{The loss curve $\bar\Lambda(E)$ in two-dimensional disk with harmonic potential $V=\frac{\omega^2}{2}|x|^2$ and $\omega=10$ with the first computed eigenfunction.}
\label{pot1.fig}
\end{figure}

\begin{figure}[tb]
\centering
\includegraphics[width=.9\columnwidth]{loss_p=2_n=2_V=harmonic_omega=1_disco.png}
\includegraphics[width=.45\columnwidth]{egen_p=2_n=2_V=harmonic_omega=1_disco_1.png}
\includegraphics[width=.45\columnwidth]{egen_p=2_n=2_V=harmonic_omega=1_disco_1_3d.png}
\caption{The loss curve $\bar\Lambda(E)$ in two-dimensional disk with harmonic potential $V=\frac{\omega^2}{2}|x|^2$ and $\omega=1$ with the first computed eigenfunction.}
\label{pot2.fig}
\end{figure}



\section{Nonlinear problems}

Our method can also be extended to nonlinear eigenvalue problems. As an illustrative example, we consider the eigenvalue problem for the $p-$Laplacian in the two-dimensional unit disk:
\begin{equation}\label{p-eigen}
\begin{cases}
-\nabla \cdot (|\nabla u|^{p-2} \nabla u) = E |u|^{p-2} u & \text{in } D, \\
u = 0 & \text{on } \partial D,
\end{cases}
\end{equation}
where $D \subset \R^2$ is the unit disk and $p > 1$.

To approximate the solution using our framework, we modify the loss function accordingly. For a fixed $E$, we define the following loss:
\begin{align*}
\Lambda_p(u, E) = & \int \left( \Delta_p u + E |u|^{p-2} u \right)^2 \, dx \\
&+ \mu \left( \int |u|^p\, dx - 1 \right)^2,
\end{align*}
where $\Delta_p u = \nabla \cdot (|\nabla u|^{p-2} \nabla u)$.

Here, the normalization condition is adapted to the natural scaling of the $p-$Laplacian: $\|u\|_p = 1$. As in the linear case, the output of the neural network is multiplied by a boundary factor to enforce homogeneous Dirichlet boundary conditions.

Since the nonlinear problem lacks an orthogonality structure among eigenfunctions, traditional sequential methods based on orthogonal projections are not directly applicable. Our approach, however, remains valid and does not rely on any such structure, which makes it suitable for problems like \eqref{p-eigen}. As a result, we are still able to identify the first two eigenvalues as the values of $E$ for which the loss function attains a sufficiently small minimum.


%\textbf{Example 3.} p-Laplacian in the unit disk
%
%\medskip
%\underline{Parameters}: \\
%$\Omega= B(0,1)$\\
%V=0\\
%Sample points: 5000\\
%mu=4000\\
%n\_epochs=2000\\
%early\_stopping\_patience=300\\
%E\_min=0\\
%E\_max=45\\
%E\_step=100
%
%\medskip
%
%Numerical computation performed in \url{https://arxiv.org/pdf/1106.0602}, Table 1.
%The authors use a variant of the Constrained Steepest
%Descent Method (CDM) to find the first eigenpair, and the Constrained Mountain Pass Algorithm
%(CMPA) to find the second eigenpair.

 
In Table \ref{table.5}, for $p=2.2$, we compare our computations with the numerical values of \eqref{p-eigen} reported in \cite{HJ11}, where a variant of the Constrained Steepest Descent Method was used to compute the first eigenpair and the Constrained Mountain Pass Algorithm to compute the second eigenpair.  Figure~\ref{loss.plap.disk.fig} shows the loss curve and the corresponding eigenfunction.

We emphasize that our method extends naturally to the computation of eigenvalues of the $p-$Laplacian in higher dimensions. in higher dimensions. To the best of our knowledge, no previous numerical results are available for this setting. Figure~\ref{loss.plap.disk.R3.fig} shows the loss curve obtained for the $p-$Laplacian on the unit ball in dimension $d=3$.


Table \ref{table.6} exhibits the same comparison when the domain $D$ is the square $[0,2]\times [0,2]$.

\begin{table*}[h!]
\centering
 
\begin{tabular}{c| c|c|c|c|c|c}
k \,& Numerical $E_{k}$ from \cite{HJ11} & Our result & Relative error (\%)&  Residual \\ 
\hline
1  & 6.5320    & 6.679258 & $2.25\times 10^{-2}$& 1.81$\times 10^0$ \\
2  & 18.395    & 19.441806 & $5.69\times 10^{-2}$&9.65$\times 10^{-1}$\\
2  & -    & 37.621714 & & 1.82$\times 10^0$
\end{tabular}
\caption{Eigenvalues of the $p-$Laplacian with $p=2.2$ in the planar unit ball  compared with numerical values obtained in \cite{HJ11}.}
\label{table.5}
 
\end{table*}


%
%\textbf{Example 4.} p-Laplacian with a potential in a square
%
%\medskip
%\underline{Parameters}: \\
%$\Omega=[0,2]\times[0,2]$\\
%V=0\\
%Sample points: 5000\\
%mu=$100+\lambda*\lambda$\\
%n\_epochs=2500\\
%early\_stopping\_patience=300\\
%E\_min=0\\
%E\_max=20\\
%E\_step=80
%
%\medskip
%
%Numerical computation performed in \url{https://arxiv.org/pdf/1106.0602}, Table 1.
%The authors use a variant of the Constrained Steepest
%Descent Method (CDM) to find the first eigenpair, and the Constrained Mountain Pass Algorithm
%(CMPA) to find the second eigenpair.

\begin{table*}[h!]
\centering
 
\begin{tabular}{c|c|c|c|c|c}
k \,& Numerical $E_{k}$ from \cite{HJ11} & Our result & Relative error (\%)&  Residual \\ 
\hline
1  & 5.4952    & 5.861383 & $1.31\times 10^{-1}$& 6.22$\times 10^{-2}$ \\
2  & 15.144    & 15.144490 & $3.23\times 10^{-5} $&1.64$\times 10^{0}$\\
\end{tabular}
\caption{Eigenvalues of the $p-$Laplacian with $p=2.2$ in the planar square of lenght 2 compared with numerical values obtained in \cite{HJ11}.}
\label{table.6}
 
\end{table*}

\normalcolor



\begin{figure}[tb]
\centering
\includegraphics[width=.9\columnwidth]{loss_n=2_p=2.2_V=0_mu=150_+_lam_x_lam_disco.png}
\includegraphics[width=.45\columnwidth]{eigen-plap_p=2.2_1.png}
\includegraphics[width=.45\columnwidth]{eigen-plap_p=2.2_1_3d.png}
\includegraphics[width=.45\columnwidth]{eigen-plap_p=2.2_2.png}
\includegraphics[width=.45\columnwidth]{eigen-plap_p=2.2_2_3d.png}
\caption{The loss curve $\bar\Lambda(E)$ in two-dimensional space for the p-Laplacian with p=2.2 in the unit disk, and the first two computed eigenfunctions.}
\label{loss.plap.disk.fig}
\end{figure}

\begin{figure}[tb]
\centering
\includegraphics[width=.7\columnwidth]{p-lap-p2.5-n3-unitball.png}
\caption{The loss curve $\bar\Lambda(E)$ in three-dimensional space for the p-Laplacian with p=2.5 in the unit disk. This gives $E_1=17.4513$ and $E_2=45.8908$.}
\label{loss.plap.disk.R3.fig}
\end{figure}













\begin{figure}[tb]
\centering
\includegraphics[width=.7\columnwidth]{well-1.png}
\includegraphics[width=.7\columnwidth]{well-1-plot.png}
\caption{The loss curve $\bar\Lambda(E)$ in three-dimensional space for the Laplacian with a double-well Gaussian potential in the unit disk.}
\label{fig.well-two}
\end{figure}



\begin{figure}[tb]
\centering
\includegraphics[width=.7\columnwidth]{well-2.png}
\includegraphics[width=.7\columnwidth]{well-2-plot.png}
\caption{The loss curve $\bar\Lambda(E)$ in three-dimensional space for the Laplacian with a Compactly supported four-well potential in the unit disk.}
\label{fig.well-four}
\end{figure}

\section{Code availability}

The Python code used in this work is openly accessible at the following URL:

\begin{center}
\url{https://github.com/amsalort/PINN-eigenvalues}
\end{center}

\section{Conclusions}

The method proposed in this work provides a flexible and effective framework for computing eigenvalues and eigenfunctions of differential operators using neural networks. Unlike classical approaches, it does not rely on the Rayleigh quotient or require orthogonality conditions between eigenfunctions. As a result, it is particularly well-suited for computing higher eigenvalues and for addressing nonlinear problems where such structures are no longer available.

Moreover, the approach naturally extends to irregular domains and high-dimensional settings, where traditional mesh-based methods become increasingly difficult to apply. Its non-sequential nature allows for the identification of multiple eigenvalues without the need to compute the entire spectrum in order. These features make it a robust and broadly applicable alternative for spectral problems in both linear and nonlinear contexts.


\section{acknowledgments}

This work was partially supported by ANPCyT under grant PICT 2019-3837 and by CONICET under grant PIP 11220150100032CO.

JFB thanks Pablo Groisman for encouraging him to teach a course on the Mathematical Foundations of Deep Learning in the fall semester of 2025, which inspired the author to pursue this problem.

\bibliography{biblio}

\appendix

\section{The minimization problem}

In this appendix, we show that the minimization problem leading to the loss curve $\Lambda(E)$ defined in \eqref{loss-curve} admits a solution. In order to keep things simple, we consider the case in which $V(x)$ is a confining potential of the form \eqref{conf.V} with $\Omega\subset \R^n$ bounded.

To this end, we must specify the function space over which the infimum is taken. The natural choice is the Sobolev space $H^2(\Omega) \cap H^1_0(\Omega)$. The main result of this section reads as follows:

\begin{teo}\label{teo.min}
Let $V(x)$ be as in \eqref{conf.V} with $\Omega\subset \R^n$ bounded. Given $E\ge 0$, there exists $u_E\in H^2(\Omega)\cap H^1_0(\Omega)$ such that
$$
\Lambda(E) = \Lambda(u_E, E) = \inf_{u\in H^2(\Omega)\cap H^1_0(\Omega)}\Lambda(u, E),
$$
where $\Lambda(u, E)$ is given in \eqref{Lambda}.
\end{teo}

\begin{proof}
The result follows from the \emph{direct method in the Calculus of Variations}. Fix $E \ge 0$, and consider a minimizing sequence $\{u_j\}_{j \in \N} \subset H^2(\Omega) \cap H^1_0(\Omega)$ such that $\|u_j\|_2 = 1$ and
$$
\Lambda(E) = \lim_{j \to \infty} \Lambda(u_j, E).
$$
In particular, the sequence $\Lambda(u_j, E)$ is bounded and nonnegative:
\begin{equation}\label{cota.Lambda}
0 \le \Lambda(u_j, E) \le C,
\end{equation}
for some constant $C > 0$.

Now observe that
\begin{align*}
\Lambda(u_j, E) &= \int_\Omega (\Delta u_j + Eu_j)^2\, dx \\
&= \int_\Omega (\Delta u_j)^2\, dx + 2E\int_\Omega u_j \Delta u_j \, dx+ E^2.
\end{align*}
Using the inequality $ab\le \epsilon a^2 + \frac{b^2}{4\epsilon}$ for any $a, b>0$ and $\epsilon >0$we estimate:
$$
\int_\Omega |u_j| |\Delta u_j|\, dx \le \epsilon \int_\Omega (\Delta u_j)^2\, dx + \frac{1}{4\epsilon}\int_\Omega u_j^2\, dx.
$$
Taking $\epsilon=1/4E$, we obtain
\begin{equation}\label{cota.Lambda2}
\Lambda(u_j, E)\ge \frac12 \int_\Omega (\Delta u_j)^2\, dx + E^2-E.
\end{equation}
Combining \eqref{cota.Lambda}, \eqref{cota.Lambda2} and the normalization $\|u_j\|_2=1$ we conclude that the sequence $\{u_j\}_{j\in\N}$ is bounded in $H^2(\Omega)\cap H^1_0(\Omega)$.

Therefore, up to a subsequence, we have $u_j \rightharpoonup u_E$ weakly in $H^2(\Omega) \cap H^1_0(\Omega)$ for some $u_E$ in that space.

Since the functional 
$$
u \mapsto \int_\Omega (\Delta u + E u)^2\, dx
$$ 
is convex in $u$, it is weakly lower semicontinuous. Thus,
$$
\int_\Omega (\Delta u_E + E u_E)^2\, dx \le \liminf_{j \to \infty} \int_\Omega (\Delta u_j + E u_j)^2\, dx.
$$
Moreover, the immersion $H^2(\Omega)\cap H^1_0(\Omega)\hookrightarrow L^2(\Omega)$ is compact, so
$$
\int_\Omega u_E^2\, dx = \lim_{j\to\infty} \int_\Omega u_j^2\, dx = 1.
$$
Hence, $u_E$ is admissible and attains the minimum.
\end{proof}
\end{document}
